\section{Related Work}

% \textbf{Image Question Answering.} Visual Question Answering \cite{VQA} requires a model to answer visual related questions via understanding both visual content and question semantics. Existing benchmarks for visual question answering \cite{goyal2017making, zhu2016visual7w, krishna2016visual} use images \cite{VQA, zhu2016visual7w, gan2017vqs} and 2types of visual comprehension questions. They have made significant progress in assessing the visual language comprehension of systems from multiple perspectives of perception. Beyond factual reasoning, more research has fallen on unusual pictures, such as the recent GPT-4 \cite{openai2023gpt4} which can interpret the meaning of funny images (An example is "Lightning Cable"),  \cite{hessel2022androids} devised benchmarks for \textit{New Yorker-caliber} cartoon. comprehension, and \cite{Whoop} which devised benchmarks for commonsense-defying images comprehension.



% Please add the following required packages to your document preamble:
% \usepackage{booktabs}
% \usepackage{multirow}
% \usepackage[table,xcdraw]{xcolor}
% If you use beamer only pass "xcolor=table" option, i.e. \documentclass[xcolor=table]{beamer}
\begin{table}[!t]
\caption{\textbf{Comparison between FunQA and other existing benchmarks.}
Compare to other datasets, FunQA revolves around the captivating realm of interesting and counter-intuitive videos. The tasks within FunQA are specifically designed to challenge the vision capabilities of models, requiring strong skills in producing an in-depth description, interpretation, and spatial-temporal reasoning.
% FunQA's domain is interesting and counter-intuitive videos, FunQA selects videos that are humanly and repeatedly screened for high quality, FunQA's tasks focus on vision, and FunQA's answers require strong description, interpretation, and spatial-temporal reasoning skills already powerful in Free-text generation. FunQA is therefore designed to look at deeper description and explanation. 
Here we clarify the abbreviation in the table. \textbf{Anno.}: Annotation; \textbf{M}: Manual, \textbf{A}: Automatic; For Input, \textbf{V, A, S,} and \textbf{B} denote Video, Audio, Subtitle, and Bounding-box; \textbf{VC} means visual-centric, \textbf{Desc.} means Description, \textbf{Expl} for Explanation, \textbf{STR} for Spatial-temporal Reasoning. For QA Tasks, \textbf{MC} denotes Multiple Choice QA, \textbf{OE} means Open Ended QA, and \textbf{FT} means Free Text QA.
\label{T:main}}
\vspace{-0.05mm}
\resizebox{1.0\linewidth}{!}{
\begin{tabular}{@{}
>{\columncolor[HTML]{FFFFFF}}l 
>{\columncolor[HTML]{FFFFFF}}l 
>{\columncolor[HTML]{FFFFFF}}c 
>{\columncolor[HTML]{FFFFFF}}c 
>{\columncolor[HTML]{FFFFFF}}c 
>{\columncolor[HTML]{FFFFFF}}c 
>{\columncolor[HTML]{FFFFFF}}c 
>{\columncolor[HTML]{FFFFFF}}c 
>{\columncolor[HTML]{FFFFFF}}c 
>{\columncolor[HTML]{FFFFFF}}c 
>{\columncolor[HTML]{FFFFFF}}c 
>{\columncolor[HTML]{FFFFFF}}c @{}}



\toprule
\cellcolor[HTML]{FFFFFF}     
& \cellcolor[HTML]{FFFFFF}                           
& \cellcolor[HTML]{FFFFFF} 
& \multicolumn{3}{c}{\cellcolor[HTML]{FFFFFF}Video} & \multicolumn{5}{c}{\cellcolor[HTML]{FFFFFF}Question Answer}                         \\ 

\cmidrule(l){4-6} \cmidrule(l){7-12} 

\multirow{-2}{*}{\cellcolor[HTML]{FFFFFF}Datasets} 
& \multirow{-2}{*}{\cellcolor[HTML]{FFFFFF}Domain}                                        
& \multirow{-2}{*}{\cellcolor[HTML]{FFFFFF}Anno.} 
& Avg length (s) & \# Clips (K) & Input                                                           
& \# QA pairs (K) & VC & Desc. & Expl. & STR & QA Task      

\\ \cmidrule(r){1-12}

MarioQA \cite{mun2017marioqa} 
& Games  & M & 3-6 & 188 & V, A & 188 & Yes & No & Yes & Yes & OE \\
TGIF-QA \cite{jang-IJCV-2019} 
& Social Media & M & 3.1 & 71.7 & V & 165.2 & Yes & Yes & No & Yes & MC \& OE \\
MovieQA \cite{tapaswi2016movieqa}                                            
& Movies & A & 202.7 & 6.77 & V, A, S & 6.4 & No & No & Yes & Yes & MC \\
CLEVRER \cite{yi2020clevrer}                                           
& Synthetic Video & M & 5 & 20 & V & 305 & Yes & No & Yes & Yes & OE \\
TVQA \cite{lei2018tvqa}    
& TV shows & M & 76.2 & 21.8 &V, A & 152.5 & No & No & Yes & No & MC \\
TVQA+ \cite{lei2019tvqa}  
& TV shows & M &7.2 & 4.2 & V, A, B & 29.4 & No & No & Yes & Yes & OE \\
Social-IQ \cite{zadeh2019social}                                          
& Web videos & M & 99 & 1.25 & V, A, S, B & 7.5 & Yes & No & Yes & No & MC \\
NExT-QA \cite{xiao2021nextqanext}                                            
& Daily life & M & 44 & 5.4 & V, A & 52 & Yes & Yes & Yes & Yes & MC \& OE \\
KnowIT-VQA \cite{garcia2020knowit}                                        
& TV shows & M & 60 & 12 & V, A, S & 24 & Yes & No & Yes & Yes & MC \\
AGQA \cite{grundemclaughlin2021agqa}                                               
& Social Media & A & 30 & 9.6 & V, A & 192 & Yes & No & Yes & Yes & OE \\
AVQA \cite{yang2022avqa}                                               
& Social Media & M & 60 & 9.3 & V, A & 57.3 & Yes & No & Yes & Yes & MC \\
STAR \cite{wu2021star}                                               
& Daily life & A & - & 22 & V & 60 & Yes & No & No & Yes & MC  \\
Env-QA \cite{Gao_2021_ICCV}                                             
& Egocentric Video & M & 20 & 23.3 & V & 85.1 & Yes & No & No & Yes & MC \\
FIBER \cite{castro2022fiber}                                              
& Daily life & M & 10 & 28 & V, A, S & 2 & No & Yes & No & Yes & OE \\
\midrule

HumorQA                                           
& Daily life & M & 7 & 1.8 & V, A & 141.3 & Yes & Yes & Yes & Yes & OE \& FT \\

CreativeQA                                     
& Performance & M & 48 & 0.9 & V, A & 78.7 & Yes & Yes & Yes & Yes & OE \& FT \\

MagicQA                                  
& Magic & M & 10 & 1.6 & V, A & 91.9 & Yes & Yes & Yes & Yes & OE \& FT \\
\midrule
\textbf{FunQA} 
& \textbf{Surprising videos} 
& M & 19 & 4.3 & V, A & \textbf{312} & Yes & Yes & Yes & Yes & \textbf{OE \& FT} \\ \bottomrule 
\end{tabular}}
\vspace{-0.5cm}
\end{table}

\noindent\textbf{Video Question Answering Benchmarks}\quad
While the visual question answering (VQA) task focuses on enhancing models' ability in image comprehension \cite{goyal2017making, zhu2016visual7w, krishna2016visual}, video question answering (VideoQA) shifts the attention towards video comprehension. VideoQA is generally more challenging than VQA as it requires a comprehensive understanding of visual content, utilization of temporal and spatial information, and exploration of relationships between recognized objects and activities~\cite{yi2020clevrer}. To address the VideoQA task, the research community has introduced various benchmarks. As depicted in Table \ref{T:main}, Most commonly used VideoQA datasets are sourced from human-centric videos like movies \cite{tapaswi2016movieqa}, TV shows \cite{lei2018tvqa, lei2019tvqa, garcia2020knowit}, and social media \cite{jang-IJCV-2019, zadeh2019social, xiao2021nextqanext, grundemclaughlin2021agqa, yang2022avqa, wu2021star, castro2022fiber}, and there are also object-centric datasets of game videos \cite{mun2017marioqa}, synthetic videos \cite{yi2020clevrer} and egocentric videos~\cite{Gao_2021_ICCV}. MovieQA \cite{tapaswi2016movieqa} and TVQA \cite{lei2018tvqa} are commonly employed by VideoQA methods, which put forward tasks related to temporal and causal reasoning. However, they rely heavily on dialogue comprehension and textual plot summaries, which severely limits the challenge of visual reasoning. TGIF-QA \cite{jang-IJCV-2019} uses animated GIFs to challenge spatial-temporal reasoning, but as most GIFs are short videos of 3 seconds, and its tasks mainly focus on action description, TGIF-QA lacks complex reasoning evaluation ability. When most datasets use multiple choice questions as QA tasks, some methods, such as NExT-QA \cite{xiao2021nextqanext}, try to join open-ended questions. NExT-QA mainly focuses on daily life videos, but the open-ended answers are mostly simple sentences containing only a few words. To sum up, most existing methods focus on ordinary videos, lack of understanding of intriguing or unexpected videos, and advanced reasoning tasks such as generating complete explanatory texts of videos remain to be explored.\\
%As depicted in Table \ref{T:main}, there is a growing interest in utilizing question responses to measure video reasoning ability \cite{tapaswi2016movieqa, lei2018tvqa, lei2019tvqa, mun2017marioqa, jang-IJCV-2019, yi2020clevrer, xiao2021nextqanext, garcia2020knowit, grundemclaughlin2021agqa, wu2021star, Gao_2021_ICCV, castro2022fiber}. However, many well-known benchmarks generate data semi-automatically from popular TV shows or films \cite{tapaswi2016movieqa, lei2018tvqa, lei2019tvqa}, resulting in tasks and answers that rely heavily on dialogue and plot summaries rather than the visual content of the video \cite{tapaswi2016movieqa, lei2018tvqa,lei2019tvqa, kim2017deepstory, zadeh2019social}. This reliance on dialogue reduces the validity of the benchmarks for measuring visual reasoning \cite{tapaswi2016movieqa, lei2018tvqa, lei2019tvqa}.\\
% \textbf{Video Question Answering Benchmarks.} 
% As shown by the benchmarks in Table \ref{T:main}, there is growing interest in using question responses to measure video reasoning ability \cite{tapaswi2016movieqa, lei2018tvqa, lei2019tvqa, mun2017marioqa, jang-IJCV-2019, yi2020clevrer, xiao2021nextqanext, garcia2020knowit, grundemclaughlin2021agqa, wu2021star, Gao_2021_ICCV, castro2022fiber}. However, the data sources for these well-known benchmarks are often generated semi-automatically from well-known TV shows or films \cite{tapaswi2016movieqa, lei2018tvqa, lei2019tvqa}, so their tasks and answers often rely on dialogue and plot summaries rather than the visual content of the video \cite{tapaswi2016movieqa, lei2018tvqa, lei2019tvqa, kim2017deepstory, zadeh2019social}, leading to models that rely more on dialogue than on visual input, thus reducing the validity of the benchmarks for measuring visual reasoning \cite{tapaswi2016movieqa, lei2018tvqa, lei2019tvqa}.
% An increasing number of benchmarks are now often designed specifically to surprise by focusing on enhancing a particular domain of the model. \cite{yang2022avqa} focuses on improving the model's ability to process audio during problem solving. Moreover, \cite{xiao2021nextqanext, grundemclaughlin2021agqa, wu2021star}  currently focus on improving spatio-temporal reasoning.
\noindent\textbf{Video Question Answering Solutions}\quad
Earlier studies have explored various models, including LSTMs and graph-based neural networks, to capture cross-modal information \cite{li2019beyond, zhao2018open}. With the advent of Transformers, video understanding models like ClipBERT \cite{leiless} and CoMVT \cite{seo2021look} emerged, focusing on the comprehension of specific frames within a video. Subsequent models, such as Violet \cite{fu2021violet}, extended their ability to encompass temporal and spatial information. However, these methods have primarily been applied to short videos. In the realm of long videos, MIST \cite{gao2022mist} stands out by achieving state-of-the-art performance and excelling in terms of computation efficiency and interpretability. Furthermore, recent Vision Language Models (VLMs) such as \cite{li2023otter,2023videochat,maaz2023videochatgpt} have showcased remarkable video understanding capabilities.\\
\noindent\textbf{Counter-Intuitive Benchmarks}\quad
While many current computer vision benchmarks primarily focus on understanding commonsense content, there is a growing interest in addressing the realm of counter-intuitiveness. Several emerging benchmarks and models cater to this domain, such as Whoops~\cite{Whoop}, which emphasizes weird, unusual, and uncanny images, and MemeGraphs \cite{kougia2023memegraphs}, which revolves around memes featuring humor and sarcasm. Furthermore, some work even challenges models to comprehend complex multimodal humor in comics~\cite{hessel2022androids}. In the realm of large vision-language models, exemplified by GPT-4 \cite{openai2023gpt4}, there is a particular focus on showcasing their ability to provide explanations for funny pictures. However, when it comes to videos, existing datasets exploring humor in TV shows or comedy tend to heavily rely on audio and narrative cues \cite{mustard, hasan-etal-2019-ur,Patro_2021_WACV}, with visual clues playing a comparatively lesser role.
% However, there is no VideoQA benchmark involving counter-intuitive videos with deeper spatial-temporal reasoning.


% \textbf{Natural Text Generation} Natural language generation(NLG) is the task to automatically generate understandable texts, typically using a non-linguistic or textual representation of information as input \cite{covington2001building}. The first generation of metrics relied on handcrafted rules that measure the surface similarity between the sentences. To illustrate, BLEU \cite{papineni2002bleu} and ROUGE \cite{lin2004rouge}, two popular metrics, rely on N-gram overlap. Current approaches largely fall into two categories. Fully learned metrics, such as CIDEr \cite{vedantam2015cider}, BEER \cite{stanojevic2014beer}, RUSE \cite{shimanaka2018ruse}, and ESIM \cite{rebecq2018esim} are trained end-to-end, and they typically rely on handcrafted features and/or learned embeddings. Conversely, hybrid metrics,such as YiSi and BERTscore \cite{zhang2020bertscore} combine trained elements. Recently, there are four commonly used NLG benchmarks: all four systems employ the reference-based COMET \cite{rei2020comet},  BLEURT \cite{sellam2020bleurt}, and BERTSCORE \cite{zhang2020bertscore} metrics, in addition to the reference-free COMET-QE \cite{chimoto-bassett-2022-comet} metric.
% Recently, an increasing number of VQA tasks have now moved towards the use of open-ended answers (see in Table \ref{T:main}), which means that answers are no longer limited to predefined options. However, most answers are simply a few words or a fixed sentence. Few of the NLG that were generated in Large Language models were used as standard answers. \cite{hessel2022androids, Whoop} sets the standard answer to the QA task to a large section of natural language generated by LLM. In order to distinguish our approach from the prevalent practice of still confining answers within a certain range in open-ended VQA tasks, we have adopted a format where our annotated answers consist of lengthy paragraphs, aiming to provide as much coherent context as possible. We refer to this answer format as Free-text. 



\begin{comment}
Producing a good VQA dataset is complex, difficult and expensive because summarizing information from videos and giving inferred question-answer pairs is laborious, and most datasets will be produced with information from other modalities, such as subtitles and audio. \\
\textbf{General VQA datasets (factoid \& inference).}\quad Most of the early factoid VQA datasets were automatically generated by algorithms and were very homogeneous in form, which makes the learned
model easy to over-fit. They also in small dataset sizes and the length of videos are short, containing only a clip of the full video makes it difficult for the questions to involve complex reasoning.\cite{yu2019activitynet} For example, in VideoQA (FIB) \cite{zhu2017uncovering}, the main task is to describe the current clip and the questions are in the form of fill-in-the-blank questions, VideoQA \cite{zeng2017leveraging} uses a question generator to generate QA pairs from description of clips of movies, moreover, MSVD-QA, MSRVTT-Q \cite{xu2017video} and YouTube2Text-QA \cite{zhao2017video} generate questions of fixed categories in different ways, and their answers are all words.
Most inference VQA datasets ---xbz---. \\
So after that, some better datasets fall into several categories, including extremely large datasets, manually filtered datasets, and datasets for a particular domain.\\
\textbf{Large-scale VQA datasets.}\quad Among them, the availability of large datasets, for example, HowTo100M \cite{miech2019howto100m} and HowToVQA69M \cite{yang2021just} crawl Youtube videos in bulk and generate a large number of question-answer pairs using subtitles, Youtube-Temporal-180M \cite{zellers2021merlot}, a dataset for learning multimodal script knowledge, derived from 6 million public YouTube videos. All these databases provides the conditions for pre-training video understanding models.\\
\textbf{Manually annotated VQA datasets.}\quad In contrast, manually filtered and manually annotated datasets provide high quality datasets along with more difficult and complex tasks, providing the data base for training and evaluating models. For example, iVQA \cite{yang2021just}, which was manually selected from HowToVQA69M, eliminated QApairs that could be answered directly from the questions, and deliberately selected difficult questions, also ActivityNet-QA \cite{yu2019activitynet}, which is dedicated to generating better QApairs, fully annotated by humans, with more complex and inferential questions involving four categories: Motion, Spatial Relationship, Temporal Relationship and free.\\
\textbf{Specific VQA Datasets.}\quad Some datasets for specific types of videos provide help to solve specific tasks in a certain domain, such as the game dataset MarioQA \cite{mun2017marioqa}, the first-person dataset EgoVQA \cite{fan2019egovqa}, the wilderness dataset WildQA \cite{castro2022wild}, the chinese VQA dataset CREATE-210K and CREATE-3M/10M \cite{zhang2022create}.
\end{comment}

% \textbf{LLM based Models.} After the remarkable results of ChatGPT [20], GPT-4 [19] and other large language models [33, 32, 9], some studies have tried to include information from other modalities in the pre-trained language models.
\begin{comment}
\textbf{Transformer-based Works.}\quad Early transformer based video understanding models such as ClipBERT \cite{leiless} and CoMVT \cite{seo2021look} target the understanding of certain frames in the video, while later models such as Violet \cite{fu2021violet} add the understanding of temporal and spatial information.
\textbf{VLM.}\quad
Recently, large-scale models have been extensively studied. As a consequence, building on top of a powerful pretrained text-only language
model has been explored in concurrent work. VisualGPT (Chen et al., 2021) showed the benefit of
initializing the weights of a VLM with a pretrained language-only model for data-efficient training.
Tsimpoukelli et al. (2021) pushed this idea even further by freezing the pretrained LM weights. The
conditioning of the model on vision is similar to the idea of prefix tuning (Lester et al., 2021; Li and
Liang, 2021; Zhou et al., 2021; Zhu et al., 2021) (also called prompt tuning) where the prefix is
encoded by a trainable visual encoder. The benefit of building on top of a strong frozen language
model is that it may enable the visual language model to retain similar powerful language-only
abilities such as few-shot language adaptation, external knowledge retrieval, or dialogue capabilities.
Several works followed this idea with architectural differences in the conditioning of the frozen
language model. For instance, MAGMA (Eichenberg et al., 2021) adds bottleneck adapters (Houlsby
et al., 2019; Sung et al., 2021) within the frozen language model; ClipCap (Mokady et al., 2021)
proposes to use a vision-to-prefix transformer to map the vision features into a prefix instead of using
a simple linear layer mapping. VC-GPT (Luo et al., 2022) moves away from the visual prefix tuning
approach. They instead explore conditioning the frozen language model by grafting new learnt 
to the frozen language model. Finally, PICA (Yang et al., 2021) and Socratic Models (Zeng et al.,
2022) propose to use off-the-shelf vision-language models (Radford et al., 2021; Zhang et al., 2021)
to communicate the content of images using language descriptions to GPT-3 (Brown et al., 2020).
% new 3
Serveral recent models such as Flamingo \cite{alayrac2022flamingo}, LLaMA \cite{zhang2023llama} and AutoAD \cite{han2023autoad} are built based on large models such as Clip, GPT, etc., which demonstrate good results in the understanding of overall video information and adaptability to different complex problems. \\ 


\textbf{Fin-tune/pretrain on VQA models use special datasets.}\quad
Several recent methods pretrain multi-modal vision-language model. For example, --- \cite{} pretrain models using large amounts of narrated videos. In contrast to task-agnostic pretraining in the previous work, we show the benefits of task-specific pretraining for our target VideoQA task.
\end{comment}

% \textbf{Humor.} Attardo \cite{attardo2010linguistic} highlights three theories of humor: Hostility \cite{goldstein1976theoretical}, Release \cite{freud1960jokes,sep-humor}, and Incongruity.

% We focus on Incongruity that people laugh at the juxtaposition of incompatible concepts and at defiance of their expectations, that is, at the incongruity between expectations and reality \cite{kant1987critique,sep-humor}, and detecting dissonance between reality and expectations can place demands on the model's temporal capabilities. There has been some debate \cite{latta1999basic,boyd2004laughter,koestler2020act} attempting to clarify the roles of juxtaposition and shifting in humor. By designing humorous QA tasks, the model needs to detect juxtaposed and transformed events, corresponding to the two challenges of Multi-grained Visual Concepts and Multi-event proposed in MIST \cite{gao2022mist}, which also requires the use of spatial information

% \textbf{Creativity.} The standard definition is bipartite \cite{runco2012standard}: Creativity requires both originality and effectiveness. Mel Rhodes's answer to the question, "What is creativity?", is this: "The word creativity is a noun naming the phenomenon in which a person communicates a new concept." This paper is more focos on this originality and novelty than with effectiveness.

% \textbf{Magic.} The art of magic consists in creating illusions of the impossible \cite{nelms2012magic}, which contains a number of conjuring effects, such as appearance, vanish, and traspostition, these effects will be achieved using techniques, props or video clips. There are specific methods theory  
% \cite{lamont2005magic} to achieve each of conjuring effects, making seemingly counter-intuitive content interpretable.

% Incongruity in Humor, Originality and Novelty in Creativity and illusions of impossible in Magic are all \textbf{counter-intuitive} phenomena that require very strong priori-knowledge and spatio-temporal reasoning ability.